\documentclass[11pt,a4paper]{article}
\usepackage{times}
\usepackage{hyperref}
\usepackage{booktabs}
\usepackage{graphicx}
\title{TritMorphLLM: Hybrid Morphological Tokenization with Explicit Composition Layers and Ternary Weights}
\author{GhostWriter42}
\date{\today}

\begin{document}
\maketitle

\begin{abstract}
We introduce TritMorphLLM, a small language model that combines a morphology-aware hybrid tokenizer with an explicit learned composition (interconnection) layer placed directly after token embeddings, and ternary (1.58-bit) weights throughout. On WikiText-103 after 10k training steps, the model achieves perfect (1.000) zero-shot morphological generalization on a systematic benchmark of 128+ unseen compounds while remaining competitive in perplexity with a vanilla BPE baseline. The explicit interconnection layer is essentially free thanks to ternary weights. Code, checkpoints, and full reproduction instructions are available at \url{https://github.com/GhostWriter42/TritMorphLLM}.
\end{abstract}

\section{Results}

\begin{table}[h]
\centering
\caption{Final 10k-step results on WikiText-103 (1024-dim, 12-layer models).}
\begin{tabular}{lrr}
\toprule
Model & Held-out Perplexity & Morphology Acc \\
\midrule
TritMorph (hybrid + composition + ternary) & 680.73 & 1.000 \\
Vanilla BPE baseline & 651.21 & 0.000 \\
\bottomrule
\end{tabular}
\end{table}

The explicit composition layer enables perfect generalization on 128+ systematically generated unseen compounds.

\section{Conclusion}
We demonstrate that a simple explicit interconnection layer on top of a hybrid tokenizer, combined with ternary weights, delivers strong morphological generalization at negligible extra cost. Future work includes larger-scale training and downstream task evaluation.

\end{document}
